% CS301 Week 8 -- Searching and elementary sorting: the search contract and
% its sorted precondition; linear search vs. binary search (halving, cost,
% variants); two applications (integer square root of n, longest common
% prefix) built by binary search on a monotone predicate; and the four
% elementary sorts the syllabus names -- insertion, selection, bubble, and
% counting sort -- with their cost analysis, stability, and in-place status.
% Difficulty: intro (course-wide default per knowledge-base-CS301.md).
% Page-grounded in Karumanchi2017 Ch.10 (Sorting, pp.504-546 PDF) and Ch.11
% (Searching, pp.547-589 PDF); Wirth2004 Sec.1.9.1 (Linear Search, p.34) and
% Sec.1.9.2 (Binary Search, pp.34-36), Sec.2.2.1/2.2.2/2.2.3 (Straight
% Insertion/Selection/Exchange, pp.46/48/50); Sedgewick2011 p.187 (the
% order-of-growth table, binary search as the logarithmic example) and
% Sec.2.1 (Elementary Sorts, p.244). Horowitz & Sahni Ch.7, Aho-Hopcroft-
% Ullman Ch.8, Neapolitan Ch.2, and Knuth Vol.3 are cited at chapter level
% only (not page-indexed) -- no page numbers are asserted for them.
% Handoff: Week 7 closed on the doubly/circular list, list-backed stack and
% queue, and polynomial manipulation, and promised "Searching -- linear versus
% binary, and the elementary sorts that put data in the order binary search
% needs." This deck pays that off. It does NOT re-teach linked-list surgery or
% the asymptotic toolkit -- Week 2's O/Theta/Omega and best/worst/average are
% reused. Continues the expression-processor thread: the calculator's symbol
% table (a list since Week 6) answers every lookup by an O(n) walk; keeping it
% sorted is what makes a binary-search lookup O(log n), and the elementary
% sorts are how the table gets sorted. Sets up Week 9 (efficient sorting and
% the comparison-sort lower bound) and Week 10 (the same table as a hash
% table).

\documentclass[aspectratio=169]{beamer}
\input{header}
\coursecode{CS301}

\title{Order Is a Superpower}
\subtitle{Week 8 --- Searching \& the Elementary Sorts}
\author[Dr. Auqib Hamid Lone]{\textbf{Dr. Auqib Hamid Lone}}
\institute[GCET Ganderbal]{PCC CS-301: Data Structures\\GCET, Ganderbal}
\date[Week 8]{\small Week 8}

\begin{document}

\frame{\titlepage}

% =============================================================================
% ACT 0 --- MOTIVATION + HANDOFF
% =============================================================================
\section{Motivation and the Handoff}

\begin{frame}{Today's Four Questions}
  \begin{enumerate}
    \item Unordered or sorted --- how much does order actually buy a search?
    \item What precondition does binary search silently assume, and what
      happens when it is violated?
    \item Can we search for a value we never stored --- the integer square
      root of $n$?
    \item Insertion, selection, bubble, counting --- which sort does each
      input reward, and at what cost?
  \end{enumerate}
  \vspace{0.2cm}
  \begin{exampleblock}{Plan}
    Write the search contract. Price linear versus binary. Search a value that
    is not in any array. Then sort: four algorithms, their traces, and their
    honest bills.
  \end{exampleblock}
\end{frame}

\begin{frame}{Handoff: What Week 7 Left Open}
  \begin{itemize}
    \item Week 7 finished the chain family: the doubly linked list
      (\texttt{p->prev}, the two-link invariant), circular rings that stop at
      \texttt{head} instead of \texttt{NULL}, the stack and queue rebuilt on
      nodes, and polynomials as ordered $(\text{coeff},\text{exp})$ lists.
    \item It closed on one promise: ``Searching --- linear versus binary, and
      the elementary sorts that put data in the order binary search needs.''
    \item The deficiency it left: every lookup so far has been a
      \emph{walk}. Our symbol table --- a list since Week 6 --- answers a
      lookup by scanning from \texttt{head}, $O(n)$ worst case.
    \item New symbols today: the search bounds \texttt{lo}, \texttt{hi},
      \texttt{mid}; the properties \key{stable} and \key{in-place}; and
      $T(n)$ for the four elementary sorts.
  \end{itemize}
  \begin{exampleblock}{This week's job}
    Show that \emph{order} turns an $O(n)$ scan into an $O(\log n)$
    probe --- then learn the cheapest ways to create that order.
  \end{exampleblock}
\end{frame}

\begin{frame}{The Calculator Looks Up an Identifier}
  \begin{itemize}
    \item The tokenizer has just produced the identifier \texttt{x}. The
      symbol table must answer: \emph{is \texttt{x} declared, and where?}
    \item Held as a linked list (Week 6), each lookup scans from
      \texttt{head} until it matches or hits \texttt{NULL}: $O(n)$
      comparisons. For a program with $2{,}000$ identifiers, that is up to
      $2{,}000$ comparisons \emph{per token}.
    \item If the same entries are kept \emph{sorted by name}, a binary search
      answers in about $\log_2 2{,}000 \approx 11$ comparisons --- roughly
      $180\times$ fewer.
    \item Two questions follow: how much does sorted order speed up search,
      and how do we get data sorted without paying more than we save?
  \end{itemize}
  \vspace{-0.1cm}
  \muted{Running application (weeks 3--10): one calculator, rebuilt each week.
    Week 8 asks what its identifier table should do to look things up fast.}
\end{frame}

\transitionslide{Every lookup so far has been a walk. Today: what does order buy?}

% =============================================================================
% ACT 1 --- LINEAR AND BINARY SEARCH
% =============================================================================
\section{Linear and Binary Search}

\sectiondivider{Section 1}{Linear and Binary Search}

\begin{frame}{Searching a Sequence: The Contract}
  \begin{itemize}
    \item Finding a name in an \emph{unordered} pile means checking one item
      at a time. Finding an entry in a \emph{sorted} directory lets you jump
      to the middle and discard half at a glance. The difference is not
      cleverness --- it is a precondition.
  \end{itemize}
  \begin{block}{Definition --- the search operation (Karumanchi Ch.~11,
    pp.547--589 PDF; Wirth Sec.~1.9.1, p.~34)
    \cite{Karumanchi2017_data_structures_algorithms_made_easy,Wirth2004_algorithms_data_structures}}
    \texttt{search(A, n, key)} returns an index $i$ with \texttt{A[i] ==
    key}, or $-1$ if \texttt{key} does not occur. \key{Linear search} makes no
    assumption about the array. \key{Binary search} \emph{requires} the array
    to be sorted in nondecreasing order --- this precondition is part of the
    contract, not an optional extra.
  \end{block}
  \begin{itemize}
    \item Worked example: in \texttt{[7, 3, 9, 1, 5, 8]}, searching $9$ feels
      the way linear search works --- scan left to right. In the sorted
      \texttt{[1, 3, 5, 7, 8, 9]}, the same target can be found by halving.
  \end{itemize}
\end{frame}

\begin{frame}[fragile]{Linear Search: Scan Until You Find It}
  \footnotesize
  \begin{block}{Linear search in C (Karumanchi Ch.~11, pp.547--589 PDF)}
\begin{semiverbatim}
int lsearch(int A[], int n, int key) {
  for (int i = 0; i < n; i++)
    if (A[i] == key) return i;   // hit: return the index
  return -1;                     // scanned the whole array: miss
}
\end{semiverbatim}
  \end{block}
  \begin{itemize}
    \item Every element is tested at most once. The loop touches
      \texttt{A[0]}, then \texttt{A[1]}, and so on, in index order --- no
      assumption about the values is used.
    \item Cost, by case:
      \begin{itemize}
        \item \good{best}: the first element matches --- $1$ comparison.
        \item \bad{worst}: absent, or last --- $n$ comparisons.
        \item average (a uniformly random hit): $(n+1)/2$ comparisons.
      \end{itemize}
    \item All three are $\Theta(n)$: linear search is a linear scan, and no
      ordering can change that.
  \end{itemize}
\end{frame}

\begin{frame}{Linear Search, Drawn}
  \begin{center}
    % Diagram: linear search of A = [7,3,9,1,5,8] for key = 1; three misses
    % then a hit at index 3. Coordinates: cells at x = 0,1.3,2.6,3.9,5.2,6.5
    % (y=0); index labels at y=-0.55; comparison marks at y=0.62.
    \begin{tikzpicture}[every node/.style={font=\scriptsize}]
      \node[draw, fill=white, minimum width=1.0cm, minimum height=0.7cm,
        align=center] (a0) at (0, 0) {7};
      \node[draw, fill=white, minimum width=1.0cm, minimum height=0.7cm,
        align=center] (a1) at (1.3, 0) {3};
      \node[draw, fill=white, minimum width=1.0cm, minimum height=0.7cm,
        align=center] (a2) at (2.6, 0) {9};
      \node[draw, fill=primary-gold!25, minimum width=1.0cm,
        minimum height=0.7cm, align=center] (a3) at (3.9, 0) {1};
      \node[draw, fill=white, minimum width=1.0cm, minimum height=0.7cm,
        align=center] (a4) at (5.2, 0) {5};
      \node[draw, fill=white, minimum width=1.0cm, minimum height=0.7cm,
        align=center] (a5) at (6.5, 0) {8};
      \node[text=negative] at (0, 0.62) {\tiny $\ne$ key};
      \node[text=negative] at (1.3, 0.62) {\tiny $\ne$ key};
      \node[text=negative] at (2.6, 0.62) {\tiny $\ne$ key};
      \node[text=positive] at (3.9, 0.62) {\tiny $=$ key};
      \node[text=neutral] at (0, -0.55) {\tiny 0};
      \node[text=neutral] at (1.3, -0.55) {\tiny 1};
      \node[text=neutral] at (2.6, -0.55) {\tiny 2};
      \node[text=neutral] at (3.9, -0.55) {\tiny 3};
      \node[text=neutral] at (5.2, -0.55) {\tiny 4};
      \node[text=neutral] at (6.5, -0.55) {\tiny 5};
      \draw[->, thick, color=primary-blue] (-1.3, 0) -- (-0.55, 0);
      \node[text=primary-blue] at (3.9, -1.2)
        {4 comparisons, then return $3$};
    \end{tikzpicture}
  \end{center}
  \vspace{-0.1cm}
  \begin{itemize}
    \item The walk is unavoidable without order: nothing in the array tells
      us $1$ cannot still be anywhere ahead, so each cell must be checked.
  \end{itemize}
\end{frame}

\begin{frame}{The Sorted Precondition --- and What Its Absence Costs}
  \begin{itemize}
    \item Binary search compares the target with the \emph{middle} element,
      then keeps only the half that could contain it. That logic is sound
      \emph{only if} everything left of the middle is smaller and everything
      right is larger --- exactly the sorted invariant.
    \end{itemize}
  \begin{alertblock}{Precondition (Karumanchi Ch.~11, pp.547--589 PDF)}
    Binary search on $A[0..n-1]$ requires $A[0] \le A[1] \le \dots \le
    A[n-1]$. Applied to unsorted data it does not error --- it returns a
    \emph{plausible wrong answer}, silently discarding the half that really
    held the key.
  \end{alertblock}
  \begin{itemize}
    \item Failure trace: search $1$ in the unsorted \texttt{[7, 3, 9, 1, 5,
      8]}. \texttt{mid = 2} (value $9$); since $1 < 9$ the algorithm discards
      indices $3..5$ --- including the $1$ at index $3$. It reports a miss.
    \item \key{Lesson}: the precondition is a \emph{contract}, not a hint.
      Binary search is only as correct as the order it is handed.
  \end{itemize}
\end{frame}

\begin{frame}[fragile]{Binary Search: Halve the Range}
  \footnotesize
  \begin{block}{Binary search in C (Wirth Sec.~1.9.2, pp.34--36; Karumanchi
    Ch.~11, pp.547--589 PDF)
    \cite{Wirth2004_algorithms_data_structures,Karumanchi2017_data_structures_algorithms_made_easy}}
\begin{semiverbatim}
int bsearch(int A[], int n, int key) {
  int lo = 0, hi = n - 1;               // current [lo, hi], inclusive
  while (lo <= hi) {
    int mid = lo + (hi - lo) / 2;       // avoids lo+hi overflow
    if (A[mid] == key) return mid;      // hit
    else if (A[mid] < key) lo = mid + 1; // key is in the right half
    else hi = mid - 1;                  // key is in the left half
  }
  return -1;                            // range empty: miss
}
\end{semiverbatim}
  \end{block}
  \begin{itemize}
    \item Each iteration either stops or halves the live range
      $[\texttt{lo},\texttt{hi}]$. Writing \texttt{lo + (hi - lo) / 2} rather
      than \texttt{(lo + hi) / 2} is not cosmetic: for large indices the sum
      can overflow \texttt{int} --- Week 1's caution about C arithmetic in
      disguise.
  \end{itemize}
\end{frame}

\begin{frame}{Binary Search, Drawn: The Range Shrinks}
  \begin{center}
    % Diagram: binary search for a key; the 8-slot range [0,8) is halved
    % each row. Coordinates: slot i spans [i, i+1] on x; row y-values
    % 2.4, 1.6, 0.8, 0.0; mid marked by a vertical rule.
    \begin{tikzpicture}[every node/.style={font=\scriptsize}]
      \foreach \i in {0,...,7}
        \node[text=neutral] at (\i + 0.5, 3.15) {\tiny \i};
      \draw[neutral, dashed] (0, 2.75) -- (8, 2.75);
      % row 1: full range [0,8), mid 4
      \draw[fill=primary-gold!20, draw=primary-gold]
        (0, 2.15) rectangle (8, 2.65);
      \draw[primary-blue, thick] (4, 2.15) -- (4, 2.65);
      \node[text=neutral, anchor=east] at (-0.15, 2.4) {$0..7$: mid $3$};
      % row 2: [0,4), mid 2
      \draw[fill=primary-gold!20, draw=primary-gold]
        (0, 1.35) rectangle (4, 1.85);
      \draw[primary-blue, thick] (2, 1.35) -- (2, 1.85);
      \node[text=neutral, anchor=east] at (-0.15, 1.6) {$0..3$: mid $1$};
      % row 3: [2,4), mid 3
      \draw[fill=primary-gold!20, draw=primary-gold]
        (2, 0.55) rectangle (4, 1.05);
      \draw[primary-blue, thick] (3, 0.55) -- (3, 1.05);
      \node[text=neutral, anchor=east] at (-0.15, 0.8) {$2..3$: mid $2$};
      % row 4: [3,4) found
      \draw[fill=positive!25, draw=positive]
        (3, -0.25) rectangle (4, 0.25);
      \node[text=positive, anchor=east] at (-0.15, 0.0) {found at $3$};
      \node[text=primary-blue, align=center] at (4, -0.85)
        {each step halves the range: blue rule $=$ \texttt{mid}};
    \end{tikzpicture}
  \end{center}
\end{frame}

\begin{frame}{Why Binary Search Is $O(\log n)$}
  \begin{itemize}
    \item Let $T(n)$ be the worst-case number of comparisons for a range of
      size $n$. One comparison, then a range of size at most $\lceil n/2
      \rceil$ --- so
      \[ T(n) = T(\lceil n/2 \rceil) + 1, \qquad T(1) = 1. \]
    \item Unrolling: after $k$ halvings the range is $n/2^{k}$, which reaches
      $1$ when $k = \log_2 n$. Hence $T(n) = \Theta(\log n)$ worst and average
      case, and $O(1)$ best case (the target is the middle at once).
  \end{itemize}
  \begin{center}
  {\footnotesize
  \begin{tabular}{lrrrr}
    \toprule
    $n$ & $10$ & $100$ & $1{,}000$ & $1{,}000{,}000$ \\
    \midrule
    comparisons $\approx \log_2 n$ & $4$ & $7$ & $10$ & $20$ \\
    \bottomrule
  \end{tabular}}
  \end{center}
  \muted{Sedgewick classifies binary search as the canonical \emph{logarithmic}
    order of growth (\cite{Sedgewick2011_algorithms}, order-of-growth table,
    p.~187; cost model, Sec.~1.4, p.~172). Doubling $n$ adds just one
    comparison.}
\end{frame}

\begin{frame}[fragile]{Binary Search Variants: First and Last Occurrence}
  \footnotesize
  \begin{block}{Lower bound --- first index with \texttt{A[i] >= key}
    (the upper bound is the same with $>$; Karumanchi Ch.~11, pp.547--589 PDF)}
\begin{semiverbatim}
int lowerBound(int A[], int n, int key) {   // first i with A[i] >= key
  int lo = 0, hi = n;                        // half-open [lo, hi)
  while (lo < hi) {
    int mid = lo + (hi - lo) / 2;
    if (A[mid] < key) lo = mid + 1;          // discard [lo, mid]
    else hi = mid;                           // keep mid as a candidate
  }
  return lo;                                 // n if every element < key
}
\end{semiverbatim}
  \end{block}
  \begin{itemize}
    \item The occurrence count is \texttt{upperBound} $-$ \texttt{lowerBound},
      in $O(\log n)$ --- no scan. Keeping \texttt{mid} as a candidate
      pinpoints the \emph{first} match.
  \end{itemize}
\end{frame}

\begin{frame}{Check Your Prediction: The Unsorted Array}
  \begin{block}{The array \texttt{[7, 3, 9, 1, 5, 8]} is unsorted. You run
    binary search for the key $3$. What does it return?}
    A: index $1$ \quad B: a wrong index or $-1$ \quad C: a runtime error
  \end{block}
  \begin{enumerate}
    \item A --- only if the search kept the half that still contains $3$; it
      does not.
    \item B --- \texttt{mid = 2} (value $9$); $3 < 9$, so it discards indices
      $3..5$ and continues left; the range collapses to $\{0\}$, the value $7$
      there is $> 3$, and the search returns $-1$ --- silently losing the $3$
      that sits at index $1$.
    \item C --- the code compiles and runs; nothing checks the precondition.
  \end{enumerate}
  \muted{Correct: B. Binary search on unsorted data never errors --- it
    returns a plausible-looking wrong answer, which is worse. Sort first.}
\end{frame}

\begin{frame}{Search Cost, Side by Side}
  \begin{center}
  {\footnotesize
  \begin{tabular}{lllll}
    \toprule
    \textbf{algorithm} & \textbf{precondition} & \textbf{best} &
      \textbf{worst} & \textbf{space} \\
    \midrule
    linear search & none & $O(1)$ & $\Theta(n)$ & $O(1)$ \\
    binary search (iterative) & sorted & $O(1)$ & $\Theta(\log n)$ & $O(1)$ \\
    \bottomrule
  \end{tabular}}
  \end{center}
  \begin{itemize}
    \item Binary search's $\Theta(\log n)$ worst case looks strictly better
      --- and it is, \emph{once you own sorted data}. The catch is the
      precondition: it costs something to establish, and it must be
      maintained after every insertion.
    \item The rest of this deck is about that cost. And before we leave
      search: two problems where binary search finds an answer that is
      \emph{not} stored anywhere in the array at all.
  \end{itemize}
\end{frame}

\transitionslide{Order halves the search. But binary search is not only for arrays of data.}

% =============================================================================
% ACT 2 --- APPLICATIONS
% =============================================================================
\section{Two Applications}

\sectiondivider{Section 2}{Two Applications of Binary Search}

\begin{frame}{Application 1: The Integer Square Root of $n$}
  \begin{itemize}
    \item Suppose we want $\lfloor \sqrt{n} \rfloor$ --- the largest integer
      $k$ with $k^2 \le n$ --- using only integer arithmetic, no floating
      point. The answer $k$ is not stored in any array.
    \item The trick: search the \emph{answer space} $0..n$ instead of an
      array. Define a predicate $P(k) = (k^2 \le n)$. As $k$ grows, $P$ is
      \emph{true} up to a point and \emph{false} after it --- a single
      true$\to$false boundary.
    \item A monotone predicate is exactly what binary search needs: find the
      largest $k$ with $P(k)$ true, and that $k$ is $\lfloor \sqrt{n}
      \rfloor$.
  \end{itemize}
  \begin{alertblock}{The pattern: binary search on a monotone predicate}
    Whenever ``is $k$ big enough?'' is a yes--no question with one flip
    point, binary search finds that point in $O(\log(\text{range}))$ probes
    --- even when no element of an array is involved.
  \end{alertblock}
\end{frame}

\begin{frame}[fragile]{Integer Square Root: Code and Predicate}
  \footnotesize
  \begin{block}{Floor of the square root by binary search on the answer
    (standard treatment; Karumanchi Ch.~11, pp.547--589 PDF)}
\begin{semiverbatim}
int isqrt(int n) {                       // floor(sqrt(n)), n >= 0
  int lo = 0, hi = n, ans = 0;
  while (lo <= hi) {
    int mid = lo + (hi - lo) / 2;
    if ((long long) mid * mid <= n) {     // P(mid) true
      ans = mid; lo = mid + 1;            // mid is feasible; push right
    } else hi = mid - 1;                  // P(mid) false; push left
  }
  return ans;                             // largest feasible k
}
\end{semiverbatim}
  \end{block}
  \begin{itemize}
    \item The \texttt{(long long)} cast widens \emph{before} the multiply:
      \texttt{mid} can reach $n$, and $n^2$ overflows \texttt{int} --- and even
      a 32-bit \texttt{long} on Windows (LLP64); Week 1's arithmetic caution
      again.
    \item Cost: the range $0..n$ halves each step: $O(\log n)$ probes,
      $O(1)$ space --- no array, no floating point.
  \end{itemize}
\end{frame}

\begin{frame}{The Square-Root Search Space, Drawn}
  \begin{center}
    % Diagram: predicate P(k) = (k*k <= 27) over k in 0..27; true for k<=5,
    % false for k>=6; answer floor(sqrt(27)) = 5. x-axis is scaled 0.2 cm per
    % unit: green [0,1.0], red (1.0,5.4]; boundary rule at x=1.0.
    \begin{tikzpicture}[every node/.style={font=\scriptsize}]
      \draw[positive, fill=positive!20] (0, 0) rectangle (1.0, 0.6);
      \draw[negative, fill=negative!15] (1.0, 0) rectangle (5.4, 0.6);
      \draw[jet, thick] (1.0, -0.15) -- (1.0, 0.75);
      \foreach \x/\lab in {0/0,1.0/5,2.0/10,3.0/15,4.0/20,5.0/25,5.4/27}
        \node[text=neutral] at (\x, -0.42) {\tiny \lab};
      \node[text=positive, anchor=east] at (-0.15, 0.3) {\tiny true};
      \node[text=negative, anchor=west] at (5.55, 0.3) {\tiny false};
      \node[text=primary-gold] at (1.0, 1.0) {$\lfloor\sqrt{27}\rfloor$};
      \node[text=primary-blue, align=center] at (2.7, 1.5)
        {one true$\to$false boundary; binary search finds it in
        $O(\log 27)$ probes};
    \end{tikzpicture}
  \end{center}
  \vspace{-0.4cm}
  \begin{center}
  {\footnotesize
  \begin{tabular}{l|cccccccc}
    \toprule
    $k$ & 0 & 1 & 2 & 3 & 4 & 5 & 6 & 7 \\
    \midrule
    $k^2$ & 0 & 1 & 4 & 9 & 16 & 25 & 36 & 49 \\
    $P(k)$ & T & T & T & T & T & T & F & F \\
    \bottomrule
  \end{tabular}}
  \end{center}
  \muted{$25 \le 27 < 36$, so the boundary sits at $k = 5$.}
\end{frame}

\begin{frame}{Application 2: The Longest Common Prefix}
  \begin{itemize}
    \item The calculator's autocomplete must find the longest prefix shared
      by all identifiers starting with a typed stem. For ``flower'',
      ``flow'', ``flight'' the shared prefix is ``fl''.
    \item Two facts make this a binary-search problem:
      \begin{enumerate}
        \item If the strings are \emph{sorted}, the longest common prefix of
          the whole set equals the common prefix of the first and the last
          string --- only two strings need comparing.
        \item The length $L$ of the common prefix of two strings is found by
          binary search: the predicate ``the first $L$ characters are equal''
          is true for small $L$ and false beyond the answer --- again a single
          true$\to$false boundary.
      \end{enumerate}
    \item So: sort the set once ($O(n \log n \cdot m)$ for strings of length
      up to $m$), then answer each prefix query in $O(m)$ using first and
      last --- no scan over all $n$ strings.
  \end{itemize}
\end{frame}

\begin{frame}[fragile]{Longest Common Prefix: Code}
  \footnotesize
  \begin{block}{LCP of two strings by binary search on the length
    (standard treatment; Karumanchi Ch.~15, pp.640--684 PDF, string search)}
\begin{semiverbatim}
int lcp2(const char *s, const char *t) {   // common prefix of s, t
  int lo = 0, hi = min(strlen(s), strlen(t)), ans = 0;
  while (lo <= hi) {
    int L = lo + (hi - lo) / 2;
    if (prefixEq(s, t, L)) { ans = L; lo = L + 1; }  // P(L) true
    else hi = L - 1;                                  // P(L) false
  }
  return ans;
}
// For a SORTED set: LCP of all = lcp2(first, last).
\end{semiverbatim}
  \end{block}
  \begin{itemize}
    \item \texttt{prefixEq(s, t, L)} costs $O(L)$; trying $O(\log m)$ lengths
      makes the query $O(m \log m)$ --- or $O(m)$ if each check is
      incremental.
    \item Sorted-set fact: LCP of all $=$ LCP(first, last) --- no $O(n)$ scan.
  \end{itemize}
\end{frame}

\begin{frame}{Longest Common Prefix, Drawn}
  \begin{center}
    % Diagram: sorted set {"flight","flow","flower"}; shared prefix "fl"
    % highlighted gold; LCP = LCP(first, last). Coordinates: letter cells of
    % width 0.55 starting at x=0; rows at y = 1.0, 0.25, -0.5.
    \begin{tikzpicture}[every node/.style={font=\scriptsize}]
      \foreach \i/\ch in {0/f,1/l,2/i,3/g,4/h,5/t}
        \node[draw, fill=primary-gold!25, minimum width=0.55cm,
          minimum height=0.55cm, inner sep=0] at (0.3+\i*0.58, 1.0)
          {\ch};
      \node[anchor=west, text=neutral] at (3.9, 1.0) {flight};
      \foreach \i/\ch in {0/f,1/l,2/o,3/w}
        \node[draw, fill=primary-gold!25, minimum width=0.55cm,
          minimum height=0.55cm, inner sep=0] at (0.3+\i*0.58, 0.25)
          {\ch};
      \node[anchor=west, text=neutral] at (2.7, 0.25) {flow};
      \foreach \i/\ch in {0/f,1/l,2/o,3/w,4/e,5/r}
        \node[draw, fill=primary-gold!25, minimum width=0.55cm,
          minimum height=0.55cm, inner sep=0] at (0.3+\i*0.58, -0.5)
          {\ch};
      \node[anchor=west, text=neutral] at (3.9, -0.5) {flower};
      \draw[positive, thick, decorate,
        decoration={brace, amplitude=4pt, mirror}]
        (0.02, -0.92) -- (1.2, -0.92);
      \node[text=positive] at (0.7, -1.4) {\tiny ``fl''};
      \node[text=primary-blue, align=center] at (5.6, 0.1)
        {sorted: LCP of all $=$\\ LCP(first, last)};
    \end{tikzpicture}
  \end{center}
  \vspace{-0.4cm}
  \begin{itemize}
    \item Checking all three strings against ``fl'' is unnecessary once they
      are sorted: only \texttt{flight} and \texttt{flower} need comparing ---
      their common prefix \emph{is} the set's.
  \end{itemize}
\end{frame}

\begin{frame}{Check Your Prediction: Prefix Length}
  \begin{block}{For the sorted set \texttt{["flight", "flow", "flower"]}, the
    predicate ``all share a prefix of length $L$'' is tested. For which $L$
    is it first \emph{false}?}
    A: $L = 1$ \quad B: $L = 2$ \quad C: $L = 3$
  \end{block}
  \begin{enumerate}
    \item A --- length $1$ is ``f'', shared by all three: true.
    \item B --- length $2$ is ``fl'', shared by all three --- still true, so
      the boundary is \emph{after} this.
    \item C --- length $3$ would require ``fli''/\allowbreak ``flo'' to agree;
      ``flight'' has ``i'' where ``flow'' has ``o''. First false at $L = 3$.
  \end{enumerate}
  \muted{Correct: C. The answer (longest common length) is $2$, the largest
    $L$ that is still true --- binary search returns $L_{\text{true}}$ and the
    prefix is its first $L$ characters.}
\end{frame}

\transitionslide{Two answers found by halving. Now: how do we create the order they need?}

% =============================================================================
% ACT 3 --- ELEMENTARY SORTING
% =============================================================================
\section{Elementary Sorting}

\sectiondivider{Section 3}{Elementary Sorting}

\begin{frame}{Why Sort at All?}
  \begin{itemize}
    \item Binary search is the reward; sorting is the price we pay for it.
      If the identifier table is sorted once and read many times, that price
      is amortized away.
    \item Sorting is not only for search. It puts ``equal'' things next to
      each other --- which is how we deduplicate, find duplicates, group
      records, and compute medians. The calculator's symbol table, its
      autocomplete set, and its report of repeated declarations all begin
      with a sort.
    \item Four algorithms, introduced not as a list to memorize but as four
      answers to ``what is the cheapest way, given \emph{this} input?'':
      insertion, selection, bubble, and counting sort.
  \end{itemize}
  \begin{exampleblock}{A look at all four, in one line each}
    Insertion: grow a sorted prefix by inserting the next element. Selection:
    repeatedly pick the minimum. Bubble: swap adjacent out-of-order pairs.
    Counting: use the key values themselves as array indices.
  \end{exampleblock}
\end{frame}

\begin{frame}{The Three Comparison Sorts, at a Glance}
  \begin{center}
  {\footnotesize
  \begin{tabular}{p{2.6cm} p{4.6cm} p{4.6cm}}
    \toprule
    \textbf{sort} & \textbf{idea} & \textbf{signature} \\
    \midrule
    insertion & keep $A[0..i-1]$ sorted; insert $A[i]$ into place &
      adaptive: near-sorted input is fast \\
    selection & find the minimum of the unsorted part; swap it forward &
      exactly $n-1$ swaps, always $\Theta(n^2)$ comparisons \\
    bubble & swap adjacent inversions; the largest ``bubbles'' to the end &
      can stop early when a pass makes no swap \\
    \bottomrule
  \end{tabular}}
  \end{center}
  \begin{itemize}
    \item All three sort \emph{by comparing pairs} of elements, so they work
      on any comparable type. Counting sort, next, breaks that property and
      gets linear time --- but only for keys from a small integer range.
    \item The first new vocabulary: a sort is \key{stable} if equal keys keep
      their original relative order, and \key{in-place} if it uses $O(1)$
      auxiliary space. Both will decide real choices.
  \end{itemize}
  \muted{Reference treatment: Knuth, \emph{TAOCP} Vol.~3
    \cite{Knuth1998_taocp_volume3}; Horowitz \& Sahni Ch.~7.}
\end{frame}

\begin{frame}[fragile]{Insertion Sort: Grow a Sorted Prefix}
  \footnotesize
  \begin{block}{Insertion sort in C (Wirth Sec.~2.2.1, p.~46; Karumanchi
    Ch.~10, pp.504--546 PDF)
    \cite{Wirth2004_algorithms_data_structures,Karumanchi2017_data_structures_algorithms_made_easy}}
\begin{semiverbatim}
void insertionSort(int A[], int n) {
  for (int i = 1; i < n; i++) {
    int key = A[i], j = i - 1;         // key is the element to place
    while (j >= 0 && A[j] > key) {     // shift larger elements right
      A[j + 1] = A[j];
      j--;
    }
    A[j + 1] = key;                    // drop key into its slot
  }
}
\end{semiverbatim}
  \end{block}
  \begin{itemize}
    \item Invariant: before iteration $i$, the subarray $A[0..i-1]$ is
      sorted. The outer loop extends it by one; the inner loop shifts until
      \texttt{key} finds its place.
    \item Worked example: sorting \texttt{[5, 2, 4, 6, 1, 3]} --- traced on
      the next slide.
  \end{itemize}
\end{frame}

\begin{frame}{Insertion Sort, Drawn}
  \begin{center}
    % Diagram: insertion sort of [5,2,4,6,1,3]; the sorted prefix (gold)
    % grows as the next key is inserted. Coordinates: 6 cells at x =
    % 0,1.0,...,5.0; rows at y = 1.4, 0.5, -0.4.
    \begin{tikzpicture}[every node/.style={font=\scriptsize}]
      % row 1 (start): all white, the first cell already a sorted prefix
      \foreach \i/\v in {0/5,1/2,2/4,3/6,4/1,5/3}
        \node[draw, fill=white, minimum width=0.9cm, minimum height=0.7cm,
          align=center] at (\i*1.0, 1.4) {\v};
      \node[draw, fill=primary-gold!35, minimum width=0.9cm,
        minimum height=0.7cm, align=center] at (0, 1.4) {5};
      \node[anchor=east, text=neutral] at (-0.6, 1.4) {start};
      % row 2: the sorted prefix is cells 0..1 (gold)
      \foreach \i/\v in {0/2,1/5,2/4,3/6,4/1,5/3}
        \node[draw, fill=white, minimum width=0.9cm, minimum height=0.7cm,
          align=center] at (\i*1.0, 0.5) {\v};
      \foreach \i/\v in {0/2,1/5}
        \node[draw, fill=primary-gold!35, minimum width=0.9cm,
          minimum height=0.7cm, align=center] at (\i*1.0, 0.5) {\v};
      \node[anchor=east, text=neutral] at (-0.6, 0.5) {insert $2$};
      % row 3: the sorted prefix is cells 0..2 (gold)
      \foreach \i/\v in {0/2,1/4,2/5,3/6,4/1,5/3}
        \node[draw, fill=white, minimum width=0.9cm, minimum height=0.7cm,
          align=center] at (\i*1.0, -0.4) {\v};
      \foreach \i/\v in {0/2,1/4,2/5}
        \node[draw, fill=primary-gold!35, minimum width=0.9cm,
          minimum height=0.7cm, align=center] at (\i*1.0, -0.4) {\v};
      \node[anchor=east, text=neutral] at (-0.6, -0.4) {insert $4$};
      \node[text=primary-gold] at (2.6, -1.15)
        {gold $=$ sorted prefix (grows by one each step)};
    \end{tikzpicture}
  \end{center}
  \vspace{-0.1cm}
  \begin{itemize}
    \item Each new key is slid left past the elements greater than it. When
      the prefix is already sorted, the inner loop runs \emph{once} --- which
      is the whole point of an adaptive sort.
  \end{itemize}
\end{frame}

\begin{frame}{Insertion Sort: The Bill}
  \begin{itemize}
    \item Count the shifts (the inner-loop body). Outer loop runs $i = 1
      \dots n-1$, so there are $n-1$ insertions.
  \end{itemize}
  \begin{center}
  {\footnotesize
  \begin{tabular}{lll}
    \toprule
    \textbf{case} & \textbf{comparisons / shifts} & \textbf{example input} \\
    \midrule
    best & $n-1$ (each key already in place) & already sorted: $O(n)$ \\
    worst & $\sum_{i=1}^{n-1} i = \dfrac{n(n-1)}{2}$ & reverse sorted: $\Theta(n^2)$ \\
    average & $\Theta(n^2)$ & random order \\
    \bottomrule
  \end{tabular}}
  \end{center}
  \begin{itemize}
    \item \good{Adaptive}: cost falls to $O(n)$ on nearly-sorted input ---
      the only one of the three comparison sorts with a linear best case.
    \item \good{Stable} (equal keys stay in order, because we shift only on
      \texttt{A[j] > key}, not $\ge$) and \good{in-place} ($O(1)$ extra
      space). These two properties are why insertion sort is the base case
      inside efficient sorts in Week 9.
  \end{itemize}
\end{frame}

\begin{frame}[fragile]{Selection Sort: Repeatedly Pick the Minimum}
  \footnotesize
  \begin{block}{Selection sort in C (Wirth Sec.~2.2.2, p.~48; Karumanchi
    Ch.~10, pp.504--546 PDF)
    \cite{Wirth2004_algorithms_data_structures,Karumanchi2017_data_structures_algorithms_made_easy}}
\begin{semiverbatim}
void selectionSort(int A[], int n) {
  for (int i = 0; i < n - 1; i++) {
    int m = i;                          // index of the minimum so far
    for (int j = i + 1; j < n; j++)
      if (A[j] < A[m]) m = j;           // find the smallest in A[i..n-1]
    swap(&A[i], &A[m]);                 // put it at the front of the rest
  }
}
\end{semiverbatim}
  \end{block}
  \begin{itemize}
    \item Invariant: after iteration $i$, the subarray $A[0..i]$ holds the
      $i+1$ smallest elements, in final position. The scan finds the next
      minimum; one swap moves it into place.
    \item Note the scan always covers the whole unsorted remainder ---
      regardless of whether the input is already sorted.
  \end{itemize}
\end{frame}

\begin{frame}{Selection Sort, Drawn}
  \begin{center}
    % Diagram: selection sort of [5,2,4,6,1,3]; green = finalised prefix;
    % gold = current minimum being selected. Cells at x = 0,1.0,...,5.0;
    % rows at y = 1.3, 0.4, -0.5.
    \begin{tikzpicture}[every node/.style={font=\scriptsize}]
      % row 1: original; the minimum 1 sits at index 4 (gold)
      \foreach \i/\v in {0/5,1/2,2/4,3/6,4/1,5/3}
        \node[draw, fill=white, minimum width=0.9cm, minimum height=0.7cm,
          align=center] at (\i*1.0, 1.3) {\v};
      \node[draw, fill=primary-gold!35, minimum width=0.9cm,
        minimum height=0.7cm, align=center] at (4, 1.3) {1};
      \node[anchor=east, text=neutral] at (-0.6, 1.3) {min is $1$};
      % row 2: after swapping 5 and 1 (index 0 green, next min 2 gold)
      \foreach \i/\v in {0/1,1/2,2/4,3/6,4/5,5/3}
        \node[draw, fill=white, minimum width=0.9cm, minimum height=0.7cm,
          align=center] at (\i*1.0, 0.4) {\v};
      \node[draw, fill=positive!25, minimum width=0.9cm,
        minimum height=0.7cm, align=center] at (0, 0.4) {1};
      \node[draw, fill=primary-gold!35, minimum width=0.9cm,
        minimum height=0.7cm, align=center] at (1, 0.4) {2};
      \node[anchor=east, text=neutral] at (-0.6, 0.4) {swap $5,1$};
      % row 3: indices 0..2 final (green)
      \foreach \i/\v in {0/1,1/2,2/3,3/6,4/5,5/4}
        \node[draw, fill=white, minimum width=0.9cm, minimum height=0.7cm,
          align=center] at (\i*1.0, -0.5) {\v};
      \foreach \i/\v in {0/1,1/2,2/3}
        \node[draw, fill=positive!25, minimum width=0.9cm,
          minimum height=0.7cm, align=center] at (\i*1.0, -0.5) {\v};
      \node[anchor=east, text=neutral] at (-0.6, -0.5) {swap $4,3$};
      \node[text=positive] at (2.6, -1.25)
        {green $=$ finalised smallest elements};
    \end{tikzpicture}
  \end{center}
  \vspace{-0.1cm}
  \begin{itemize}
    \item Row 1: the scan over the whole array finds $1$ (index $4$); swap it
      with index $0$. Row 2: the scan now finds $2$; it is already in place.
      Row 3: finds $3$; swap with $4$.
  \end{itemize}
\end{frame}

\begin{frame}{Selection Sort: The Bill (and a Warning)}
  \begin{itemize}
    \item The inner loop comparisons do not depend on the data at all:
      $\sum_{i=0}^{n-2} (n-1-i) = \dfrac{n(n-1)}{2}$ --- every input.
  \end{itemize}
  \begin{center}
  {\footnotesize
  \begin{tabular}{lll}
    \toprule
    \textbf{case} & \textbf{comparisons} & \textbf{swaps} \\
    \midrule
    best & $\Theta(n^2)$ & $0$ (already sorted) \\
    worst & $\Theta(n^2)$ & $n-1$ \\
    \bottomrule
  \end{tabular}}
  \end{center}
  \begin{itemize}
    \item \bad{Not adaptive}: a sorted input is scanned just as hard as a
      reversed one --- always $\Theta(n^2)$ comparisons. Its one virtue is
      the minimum number of \emph{writes} ($n-1$ swaps), useful when moving
      data is far costlier than comparing it.
    \item \bad{Not stable}: the long-range swap can leap an equal key over
      its twin. (Trace \texttt{[2, 2, 1]} to see the two $2$s change order.)
    \item \good{In-place}: $O(1)$ auxiliary space.
  \end{itemize}
\end{frame}

\begin{frame}[fragile]{Bubble Sort: Swap Neighbours, Bubble the Max Up}
  \footnotesize
  \begin{block}{Bubble sort with an early exit in C (Wirth Sec.~2.2.3, p.~50;
    Karumanchi Ch.~10, pp.504--546 PDF)
    \cite{Wirth2004_algorithms_data_structures,Karumanchi2017_data_structures_algorithms_made_easy}}
\begin{semiverbatim}
void bubbleSort(int A[], int n) {
  for (int i = 0; i < n - 1; i++) {
    int swapped = 0;
    for (int j = 0; j < n - 1 - i; j++)
      if (A[j] > A[j + 1]) { swap(&A[j], &A[j + 1]); swapped = 1; }
    if (!swapped) break;     // a clean pass: already sorted
  }
}
\end{semiverbatim}
  \end{block}
  \begin{itemize}
    \item Invariant: after pass $i$, the largest $i+1$ elements are final at
      the right end, so each pass can stop $i$ slots earlier.
    \item The \texttt{swapped} flag is the practical touch: a pass with no
      swap proves the array is sorted, so we stop.
  \end{itemize}
\end{frame}

\begin{frame}{Bubble Sort, Drawn}
  \begin{center}
    % Diagram: one bubble pass over [5,2,4,6,1,3]; adjacent swaps push the
    % largest to the right end. Cells at x = 0,1.0,...,5.0; two rows at
    % y = 1.2 and 0.1; swap arrows between adjacent cells in row 1.
    \begin{tikzpicture}[every node/.style={font=\scriptsize}]
      \foreach \i/\v in {0/5,1/2,2/4,3/6,4/1,5/3}
        \node[draw, fill=white, minimum width=0.9cm, minimum height=0.7cm,
          align=center] at (\i*1.0, 1.2) {\v};
      \foreach \i in {0,1,2,4}
        \draw[<->, thick, color=negative] (\i*1.0+0.45, 1.62)
          -- (\i*1.0+0.55, 1.62);
      \node[anchor=east, text=neutral] at (-0.6, 1.2) {pass 1};
      \foreach \i/\v in {0/2,1/4,2/5,3/1,4/3,5/6}
        \node[draw, fill=white, minimum width=0.9cm, minimum height=0.7cm,
          align=center] at (\i*1.0, 0.1) {\v};
      \node[draw, fill=primary-gold!25, minimum width=0.9cm,
        minimum height=0.7cm, align=center] at (5, 0.1) {6};
      \node[anchor=east, text=neutral] at (-0.6, 0.1) {result};
      \node[text=primary-gold, align=center] at (2.6, -0.7)
        {$6$ bubbles to its final slot; the next pass stops one slot sooner};
    \end{tikzpicture}
  \end{center}
  \vspace{-0.45cm}
  \begin{itemize}
    \item Comparing neighbours \texttt{(A[j], A[j+1])} and swapping when out
      of order moves the maximum rightward one step per comparison, so after
      one full pass it rests at the end.
  \end{itemize}
\end{frame}

\begin{frame}{Bubble Sort: The Bill}
  \begin{center}
  {\footnotesize
  \begin{tabular}{lll}
    \toprule
    \textbf{case} & \textbf{comparisons} & \textbf{note} \\
    \midrule
    best & $n-1$ & sorted input: pass $1$ swaps nothing, early exit \\
    worst & $\dfrac{n(n-1)}{2}$ & reverse sorted: $\Theta(n^2)$ \\
    average & $\Theta(n^2)$ & random order \\
    \bottomrule
  \end{tabular}}
  \end{center}
  \begin{itemize}
    \item \good{Adaptive} only \emph{with} the early-exit flag; the naive
      version is $\Theta(n^2)$ on every input, including sorted data.
    \item \good{Stable}: equal adjacent keys are never swapped, because we
      swap only on strict $>$ --- their relative order is preserved.
    \item \good{In-place}. Comparison count equals insertion sort's on
      average, but bubble does far more \emph{swaps}; insertion sort is the
      better practical choice among the three.
  \end{itemize}
  \muted{All three comparison sorts hit the same $\Theta(n^2)$ wall. Week 9
    asks whether any comparison sort can do better --- and shows counting
    sort, next, sidestepping comparisons entirely.}
\end{frame}

\begin{frame}{Counting Sort: Sort Without Comparing}
  \begin{itemize}
    \item All three sorts above compare keys pairwise. But if the keys are
      small non-negative integers, the key \emph{is already an array index} ---
      so it can be counted, not compared.
    \item Counting sort tallies how many times each value occurs, converts
      those counts to positions with a running total, and scatters each
      element to its slot. It never asks ``is $A[i] < A[j]$?''
    \item The precondition: keys lie in a known small range $0..k-1$. The
      trade: linear time, but $\Theta(n + k)$ extra space --- and it is only
      correct for keys that fit that range.
  \end{itemize}
  \begin{exampleblock}{Where the calculator uses it}
    Counting the frequency of each byte or token class, and bucketing
    identifiers by length, are both one-pass counting sorts --- no
    comparisons at all.
  \end{exampleblock}
\end{frame}

\begin{frame}[fragile]{Counting Sort: Code}
  \footnotesize
  \begin{block}{Counting sort with a stable scatter (Karumanchi Ch.~10,
    pp.504--546 PDF; Horowitz \& Sahni Ch.~7, standard treatment)
    \cite{Karumanchi2017_data_structures_algorithms_made_easy,HorowitzSahni2008_fundamentals_data_structures}}
\begin{semiverbatim}
void countingSort(int A[], int n, int k) {   // keys in 0..k-1
  int *count = calloc(k, sizeof(int));       // zero-initialised tally
  for (int i = 0; i < n; i++) count[A[i]]++;
  for (int v = 1; v < k; v++) count[v] += count[v - 1];  // prefix sums
  int *out = malloc(n * sizeof(int));
  for (int i = n - 1; i >= 0; i--)           // reverse pass => stable
    out[--count[A[i]]] = A[i];
  for (int i = 0; i < n; i++) A[i] = out[i];
  free(out); free(count);                    // Week 1: every malloc freed
}
\end{semiverbatim}
  \end{block}
  \begin{itemize}
    \item \texttt{calloc(k, ...)} zeroes the tally, so no counting slot is
      read before it is written. The reverse pass makes equal keys land in
      their original relative order --- that is what makes this sort stable.
  \end{itemize}
\end{frame}

\begin{frame}{Counting Sort, Drawn: From Tallies to Positions}
  \begin{center}
    % Diagram: for A = [4,2,2,5,3,3,1] (keys 0..5): tally per value, then
    % prefix sums give each value's position. Cells at x = 0,1.0,...,5.0 with
    % value labels beneath; two rows at y = 0.9 (count) and y = -0.1 (prefix).
    \begin{tikzpicture}[every node/.style={font=\scriptsize}]
      \foreach \i/\c in {0/0,1/1,2/2,3/2,4/1,5/1}
        \node[draw, fill=light-bg, minimum width=0.9cm,
          minimum height=0.7cm, align=center] (c\i) at (\i*1.0, 0.9) {\c};
      \node[anchor=east, text=neutral] at (-0.6, 0.9) {tally};
      \foreach \i/\p in {0/0,1/1,2/3,3/5,4/6,5/7}
        \node[draw, fill=primary-gold!20, minimum width=0.9cm,
          minimum height=0.7cm, align=center] at (\i*1.0, -0.1) {\p};
      \node[anchor=east, text=neutral] at (-0.6, -0.1) {prefix};
      \foreach \i in {0,...,5}
        \node[text=neutral] at (\i*1.0, -0.65) {\tiny value \i};
      \node[text=primary-blue, align=center] at (2.5, -1.35)
        {prefix sum $=$ where the next occurrence of that value goes};
    \end{tikzpicture}
  \end{center}
  \vspace{-0.1cm}
  \begin{itemize}
    \item $A = [4,2,2,5,3,3,1]$: value $2$ occurs twice, value $3$ twice. The
      prefix row says the $1$st slot belongs to value $0$, the slots
      $1$--$2$ to value $2$, and so on. Scattering in reverse fills them
      right-to-left, preserving equal keys' order.
  \end{itemize}
\end{frame}

\begin{frame}{Counting Sort: The Bill}
  \begin{center}
  {\footnotesize
  \begin{tabular}{llll}
    \toprule
    \textbf{phase} & \textbf{work} & \textbf{cost} \\
    \midrule
    tally & one pass over $A$ & $\Theta(n)$ \\
    prefix sums & one pass over $0..k-1$ & $\Theta(k)$ \\
    scatter & one reverse pass over $A$ & $\Theta(n)$ \\
    \midrule
    total & & $\Theta(n + k)$ time, $\Theta(n + k)$ space \\
    \bottomrule
  \end{tabular}}
  \end{center}
  \begin{itemize}
    \item \good{Linear time} when $k = O(n)$. But watch the range: if keys
      span a huge range (say $0..10^9$), the tally array is enormous and
      counting sort is worse than useless --- this is why it is restricted to
      small integer keys.
    \item \good{Stable} (with the reverse scatter), \bad{not in-place}
      ($\Theta(n + k)$ auxiliary). It sorts \emph{without comparing keys} ---
      so the $\Theta(n \log n)$ comparison lower bound of Week 9 does not
      apply to it.
  \end{itemize}
\end{frame}

\begin{frame}{The Four Sorts, Compared}
  \begin{center}
  {\footnotesize
  \begin{tabular}{llllll}
    \toprule
    \textbf{sort} & \textbf{best} & \textbf{worst} & \textbf{space} &
      \textbf{stable} & \textbf{in-place} \\
    \midrule
    insertion & $O(n)$ & $\Theta(n^2)$ & $O(1)$ & yes & yes \\
    selection & $\Theta(n^2)$ & $\Theta(n^2)$ & $O(1)$ & no & yes \\
    bubble & $O(n)$ & $\Theta(n^2)$ & $O(1)$ & yes & yes \\
    counting & $\Theta(n+k)$ & $\Theta(n+k)$ & $O(n+k)$ & yes & no \\
    \bottomrule
  \end{tabular}}
  \end{center}
  \begin{itemize}
    \item Read the table as a workload guide, not a ranking:
      \begin{itemize}
        \item \key{nearly sorted} input $\to$ insertion sort: its $O(n)$ best
          case is the whole point.
        \item \key{small integer keys} in a known range $\to$ counting sort:
          linear, but pay the $O(n+k)$ space.
        \item \key{minimal writes} matters more than comparisons $\to$
          selection sort: exactly $n-1$ swaps.
      \end{itemize}
    \item No single winner --- every row is right for \emph{some} input. The
      skill is stating the workload, then choosing.
  \end{itemize}
\end{frame}

\begin{frame}{Check Your Prediction: Which Sort?}
  \begin{block}{Two inputs. Which elementary sort fits each best? \quad
    (i) an array that is already almost sorted \quad (ii) $10^6$ ages in
    $0..120$}
  \end{block}
  \begin{enumerate}
    \item For (i), \key{insertion sort}: nearly every key is already in
      place, so the inner loop rarely shifts --- it runs in $O(n)$.
    \item For (ii), \key{counting sort}: keys are integers in a tiny range,
      so $k = 121 \ll n$; the sort is $\Theta(n + k) \approx \Theta(n)$.
    \item Selection or bubble would be $\Theta(n^2)$ on both --- the wrong
      tool for either input.
  \end{enumerate}
  \muted{Matching the algorithm to the input's structure is the real skill
    this week tests, not memorizing the update loops.}
\end{frame}

\transitionslide{Four sorts, one job. Now the vocabulary that decides between them.}

% =============================================================================
% ACT 4 --- PROPERTIES, REVIEW, AND CLOSE
% =============================================================================
\section{Stability, Review, and Close}

\sectiondivider{Section 4}{Stability and the Close}

\begin{frame}{Stability and In-Place, Made Concrete}
  \begin{itemize}
    \item \key{Stable}: equal keys keep their original relative order. This
      matters when records are sorted by a \emph{secondary} key after already
      being sorted by a primary one --- a stable sort preserves the earlier
      ordering within ties.
  \end{itemize}
  \begin{exampleblock}{Worked example --- stability}
    Sort \texttt{[(A,2), (B,1), (C,2)]} by the number. A stable sort yields
    \texttt{[(B,1), (A,2), (C,2)]} --- \texttt{A} stays before \texttt{C}
    (both key $2$). An unstable sort may return \texttt{[(B,1), (C,2),
    (A,2)]} --- same keys, different record order.
  \end{exampleblock}
  \begin{itemize}
    \item \key{In-place}: $O(1)$ auxiliary space beyond the input. Insertion,
      selection, and bubble sort are in-place; counting sort needs a
      $\Theta(n+k)$ output buffer and a $\Theta(k)$ tally, so it is not.
    \item \good{Stable}: insertion, bubble, counting. \bad{Not stable}:
      selection (its long-range swap can leap an equal key).
  \end{itemize}
\end{frame}

\begin{frame}{Three Classic Mistakes}
  \begin{itemize}
    \item \bad{Binary searching unsorted data} --- no crash, just a
      plausible wrong answer. The precondition is a contract: sort first, or
      use linear search.
    \item \bad{Forgetting the overflow in \texttt{mid}} ---
      \texttt{(lo + hi) / 2} overflows \texttt{int} for large indices; write
      \texttt{lo + (hi - lo) / 2}. The same trap as \texttt{mid * mid} in
      \texttt{isqrt} and \texttt{(long long)} casts in C.
    \item \bad{Claiming a sort is $\Theta(n \log n)$ ``because it sorts''} ---
      the elementary sorts are all $\Theta(n^2)$ worst case. State the
      \emph{case} and the \emph{bound}, and check whether the algorithm even
      compares keys (counting sort does not).
  \end{itemize}
  \begin{exampleblock}{How to revise}
    For each sort, write one sentence: the invariant, the best case, the
    worst case, and whether it is stable. Then trace one pass by hand. The
    test is traces and cases, not definitions.
  \end{exampleblock}
\end{frame}

\begin{frame}{Week 8 Summary}
  \begin{itemize}
    \item \key{The search contract}: linear search needs no order, $\Theta(n)$
      worst case; binary search \emph{requires} sorted input and is
      $\Theta(\log n)$ worst case --- but silently wrong on unsorted data.
    \item \key{Binary search on a monotone predicate} answers questions about
      values that are not stored: $\lfloor \sqrt{n} \rfloor$ (search the
      answer space $0..n$), and the longest common prefix (search the prefix
      length; sort once, then compare first and last).
    \item \key{Elementary sorts}: insertion (adaptive, stable, $O(n)$ best);
      selection (always $\Theta(n^2)$, $n-1$ swaps, unstable); bubble
      (stable, adaptive with an early exit); counting ($\Theta(n+k)$, stable,
      $O(n+k)$ space, small integer keys).
    \item \key{Dictionary at the end}: stable, in-place, best/worst case ---
      the three properties that decide which sort fits an input.
  \end{itemize}
  \begin{exampleblock}{Next week}
    Why can no \emph{comparison} sort beat $\Theta(n \log n)$, and which
    algorithms escape that bound? Efficient sorting: quick, merge, shell,
    radix, heap.
  \end{exampleblock}
\end{frame}

\begin{frame}{Before You Go: Predict}
  \begin{block}{Three one-line predictions}
    Answer from the algorithms, then check in the lab.
  \end{block}
  \begin{enumerate}
    \item Binary search on a sorted array of $1{,}000$ elements: about how
      many comparisons in the worst case, and how many would linear search
      need?
    \item Sorting \texttt{[3, 1, 3, 2]} --- which of insertion, selection,
      bubble, counting is stable, and what order do the three $3$s/ties take?
    \item Is the ``first occurrence'' variant of binary search different
      from ordinary binary search? Why return a candidate \texttt{mid}
      instead of discarding it?
  \end{enumerate}
  \muted{Answers: $\approx 10$ ($\log_2 1000$) versus up to $1{,}000$ ---
    order is worth roughly a hundredfold here. Insertion, bubble, and
    counting are stable; selection may reorder equal keys. Yes, it differs:
    keeping \texttt{mid} as a candidate when \texttt{A[mid] == key} lets the
    search continue left to find the \emph{first} match.}
\end{frame}

\begin{frame}[allowframebreaks]{References}
  \scriptsize
  \bibliographystyle{plain}
  \bibliography{../../Bibliography_base}
\end{frame}

\end{document}
